% Stage 8 restored fragment — 2.2.1 Time-domain evolution (archive lines 2469-2734)
\subsection{Time-domain evolution}
\label{sec:ch221}

Time-harmonic phasors match fixed-frequency range measurements; $\exp(-\j\omega t)$ fixes the sign of reactive terms \cite{stratton1941}. Maxwell's retarded export to $\Gamma_R$ remains the governing constraint \cite{maxwell1873}.
The physical problem is causal propagation: a switched feed, a pulsed radar, or a transient
immunity test excites fields whose support in $(\mathbf{r},t)$ must respect finite signal
speed \cite{maxwell1873}. Time-harmonic phasors are how network analyzers, compact
ranges, and chamber receivers report steady-state radiation at each frequency step
\cite{stratton1941}. The $\exp(-\jj\omega t)$ convention fixed in Chapter~\ref{ch:01} aligns
Poynting flux with wattmeter readings and VNA phase references \cite{poynting1884}. Phasor
analysis is best viewed as a steady-state reduction of a fundamentally causal theory rather
than as the primary formulation of electromagnetic dynamics \cite{harrington2001}. Even small
phase errors can rotate synthesized patterns when near-field data are back-propagated
\cite{hansen1988}. The time-domain picture below states the causal support that every VNA tone
must inherit when $\mathcal{L}_\omega$ is interpreted as a frequency slice rather than an
independent postulate \cite{stratton1941}.

Subsection~\ref{sec:ch221} supplies the time-domain foundation underlying every harmonic
formulation used later in the chapter. It introduces the first-order Maxwell evolution
system, the associated evolution operator, and the causal support restrictions that all
physically admissible frequency-domain solutions must inherit \cite{stratton1941}.
Time-domain evolution specifies retarded trajectories: the field observed at
$(\mathbf{r},t)$ depends only on source events contained within the corresponding past
light cone \cite{stratton1941}. Any numerical or analytical approximation that violates
this support structure introduces nonphysical precursors and loses its connection to
Maxwell causality \cite{harrington2001}. Local Maxwell balances remain those of
Eq.~\eqref{eq:ch1.ibvp-time-domain}; they are cited. Function spaces from
Section~\ref{sec:ch21} enter through time-dependent embeddings developed below; Sobolev and
trace declarations are not repeated \cite{stratton1941}.
Let \(\Omega\subset\mathbb{R}^{3}\) be a Lipschitz domain with boundary
\(\Gamma=\partial\Omega\). Real time-domain fields
\(\E(\mathbf{r},t)\), \(\H(\mathbf{r},t)\) use SI units and the same constitutive
class as \ref{ch:01}.
The evolution problem is posed on trajectories of the electromagnetic state rather
than on isolated snapshots at a single frequency.
A pulsed or switched excitation on a feed line makes the distinction visible: the field at
$(\mathbf{r},t)$ cannot respond until source events inside the backward light cone have
occurred. Time-domain evolution is therefore the laboratory picture behind every harmonic
slice $\mathcal{S}_\omega$ used later in the chapter \cite{stratton1941}. The state vector
$\mathbf{U}(t)$ packages the coupled $(\E,\H)$ fields at frozen time $t$ into a single
object on which the Maxwell generator acts; this packaging mirrors the phasor state $\mathbf{u}_\omega$ of Eq.~\eqref{eq:ch2.maxwell-state-vector},
with time retained rather than suppressed \cite{harrington2001}. Collect the
electromagnetic state at time \(t\) as
\begin{equation}
\label{eq:ch2.time-state-vector}
\mathbf{U}(t)
=
\begin{bmatrix}\E(\cdot,t)\\ \H(\cdot,t)\end{bmatrix}
\in\mathcal{V}(\Omega),
\end{equation}
where \(\mathcal{V}(\Omega)\) is the time-dependent product space of square-integrable
field pairs compatible with Eq.~\eqref{eq:ch1.admissible-space} at each frozen time
slice \cite{stratton1941}. Equation~\eqref{eq:ch2.time-state-vector} is the
time-domain analogue of Eq.~\eqref{eq:ch2.maxwell-state-vector}; phasor symbols are
reserved for the \(\exp(-\jj\omega t)\) reduction in Subsection~\ref{sec:ch222}. The time-domain Maxwell
system is the causal law from which every harmonic operator in Section~\ref{sec:ch22} is
obtained by phasor substitution. Physically, Eqs.~\eqref{eq:ch2.time-maxwell-evolution}
state that changing $\E$ in space drives $\H$ in time through Faraday's law, and changing
$\H$ in space drives $\E$ in time through Amp\`ere's law with impressed current as external
drive \cite{maxwell1873}. From an antenna engineer's perspective, this is the
system that finite-difference time-domain codes march forward in time; the harmonic solvers
of Subsection~\ref{sec:ch222} are steady-state limits of the same law when all data are
single-frequency \cite{harrington2001}. The time-domain Maxwell system imported from
Eq.~\eqref{eq:ch1.ibvp-time-domain} reads in vacuum constitutive form as
\begin{equation}
\label{eq:ch2.time-maxwell-evolution}
\begin{aligned}
\curl\E(\mathbf{r},t)&=-\mu_0\,\partial_t\H(\mathbf{r},t),\\
\curl\H(\mathbf{r},t)&=\mathbf{j}_{\mathrm{imp}}(\mathbf{r},t)+\varepsilon_0\,\partial_t\E(\mathbf{r},t),\\
\divg\E(\mathbf{r},t)&=\frac{\rho(\mathbf{r},t)}{\varepsilon_0},\qquad
\divg\H(\mathbf{r},t)=0,
\end{aligned}
\end{equation}
with impressed current density \(\mathbf{j}_{\mathrm{imp}}(\mathbf{r},t)\) and charge
density \(\rho(\mathbf{r},t)\). Equation~\eqref{eq:ch2.time-maxwell-evolution}
packages local dynamics as a first-order system in \((\E,\H)\) and \(\partial_t\);
the hyperbolic coupling is the time-domain face of the curl structure in
Eq.~\eqref{eq:ch1.maxwell-operator-block} \cite{stratton1941}. Charge continuity
requires
\begin{equation}
\label{eq:ch2.charge-continuity-time}
\divg\mathbf{j}_{\mathrm{imp}}(\mathbf{r},t)+\partial_t\rho(\mathbf{r},t)=0
\end{equation}
in the sense of distributions when \(\rho\) is present \cite{stratton1941}. Pairs
\((\mathbf{j}_{\mathrm{imp}},\rho)\) violating
Eq.~\eqref{eq:ch2.charge-continuity-time} lie outside the time-domain admissible
class parallel to Eq.~\eqref{eq:ch2.source-admissibility}.
\begin{figure}[t]
\centering
\ChOneFigBlock{%
\begin{tikzpicture}[ch1fig, node distance=5mm and 7mm]
 \node[ch1box] (ic) {Initial data\\$(\E_0,\H_0)$};
 \node[ch1box, right=7mm of ic] (src) {Impressed\\$\mathbf{j}_{\mathrm{imp}}(\cdot,t)$};
 \node[ch1box, right=7mm of src] (ev) {Evolution\\$\mathcal{T}(t)$};
 \node[ch1box, right=7mm of ev] (st) {State $\mathbf{U}(t)$};
 \draw[ch1flow] (ic.east) -- (ic.east |- ev.west);
 \draw[ch1flow] (ic.east |- ev.west) -- (ev.west);
 \draw[ch1flow] (src.east) -- (src.east |- ev.west);
 \draw[ch1flow] (src.east |- ev.west) -- (ev.west);
 \draw[ch1flow] (ev.east) -- (ev.east |- st.west);
 \draw[ch1flow] (ev.east |- st.west) -- (st.west);
\end{tikzpicture}%
}
\caption{Time-domain evolution: initial data and impressed sources drive the
semigroup \(\mathcal{T}(t)\) to field state \(\mathbf{U}(t)\) under boundary class
\(\mathcal{B}_\Gamma\).}
\label{fig:ch2.time-evolution-measurement chain}
\end{figure}
Boundary data enter through the same 
\(\mathcal{B}_\Gamma\) of Eq.~\eqref{eq:ch1.boundary-data-operator}, now evaluated
on time-dependent traces \(\gamma_\tau[\mathbf{U}(t)]\in\mathcal{T}_\Gamma\) for
each \(t\ge 0\) \cite{stratton1941}. The closed time-domain problem
is to find \(\mathbf{U}(t)\) such that
\begin{equation}
\label{eq:ch2.time-ibvp-package}
\begin{cases}
\partial_t\mathbf{U}=\mathcal{A}\,\mathbf{U}+\mathbf{F}(t), & \mathbf{r}\in\Omega,\ t>0,\\
\mathbf{U}(0)=\mathbf{U}_0, & \mathbf{r}\in\Omega,\\
\mathcal{B}_\Gamma[\mathbf{U}(t)]=\mathbf{g}_\Gamma(t), & \mathbf{r}\in\Gamma,\ t\ge 0,
\end{cases}
\end{equation}
where \(\mathcal{A}\) is the Maxwell generator induced by
Eq.~\eqref{eq:ch2.time-maxwell-evolution} on
\(\mathcal{D}(\mathcal{A})\subset\mathcal{V}(\Omega)\) and
\(\mathbf{F}(t)\) collects impressed-source rows \cite{stratton1941}. Equation
\eqref{eq:ch2.time-ibvp-package} is the time-domain counterpart of
Eq.~\eqref{eq:ch2.closed-maxwell-bvp}; uniqueness on a fixed branch is inherited
from Section~\ref{sec:ch01.1.3} when data classes coincide with the harmonic limit
\cite{stratton1941}. Define the evolution operator (semigroup) \(\mathcal{T}(t)\) by
\begin{equation}
\label{eq:ch2.time-semigroup-definition}
\mathcal{T}(t):\mathbf{U}_0\longmapsto \mathbf{U}(t),
\qquad
\mathcal{T}(0)=\mathrm{id},\qquad
\mathcal{T}(t+s)=\mathcal{T}(t)\mathcal{T}(s),
\end{equation}
when \(\mathbf{F}\equiv\mathbf{0}\) and \(\mathcal{B}_\Gamma\) is homogeneous. With
forcing,
\begin{equation}
\label{eq:ch2.time-evolution-forced}
\mathbf{U}(t)
=
\mathcal{T}(t)\mathbf{U}_0
\int_0^t \mathcal{T}(t-s)\,\mathbf{F}(s)\,ds,
\end{equation}
the Duhamel representation on \(\mathcal{D}(\mathcal{A})\)
\cite{stratton1941}. Equation~\eqref{eq:ch2.time-evolution-forced}
makes causality explicit: \(\mathbf{U}(t)\) depends only on
\(\mathbf{F}(s)\) for \(s\le t\). The time-domain \emph{solution map} sends impressed
data and initial state to trajectories:
\begin{equation}
\label{eq:ch2.time-domain-solution-map}
\mathcal{S}_{t}:\big(\mathbf{U}_0,\mathbf{j}_{\mathrm{imp}}\big)\longmapsto
\{\mathbf{U}(t)\}_{t\ge 0}.
\end{equation}
Harmonic \(\mathcal{S}_\omega\) of Eq.~\eqref{eq:ch2.maxwell-solution-map} is a
single-frequency slice of Eq.~\eqref{eq:ch2.time-domain-solution-map} when all data
are proportional to \(\exp(-\jj\omega t)\) and \(\mathcal{B}_\Gamma\) is fixed
\cite{stratton1941}; the equivalence is proved in Subsection~\ref{sec:ch223}.
Retarded representation specializes Eq.~\eqref{eq:ch2.time-evolution-forced} when
Green kernels exist:
\begin{equation}
\label{eq:ch2.retarded-convolution-form}
\E(\mathbf{r},t)
=
\E_{\mathrm{hom}}(\mathbf{r},t)
-\mu_0\,\partial_t
\int_0^t\!\!\int_\Omega
\mathcal{G}(\mathbf{r},\mathbf{r}',t-t')\cdot
\mathbf{j}_{\mathrm{imp}}(\mathbf{r}',t')\,d^{3}r'\,dt',
\end{equation}
with \(\mathcal{G}(\mathbf{r},\mathbf{r}',t-t')=\mathbf{0}\) for \(t-t'<0\)
\cite{stratton1941}. Equation~\eqref{eq:ch2.retarded-convolution-form} refines the
Eq.~\eqref{eq:ch2.causality-retarded} without constructing \(\mathcal{G}\)
here; kernel development is in Section~\ref{sec:ch23}. The support condition
\begin{equation}
\label{eq:ch2.causal-support-constraint}
\operatorname{supp}_{(\mathbf{r},t)}\E
\subseteq
\big\{(\mathbf{r},t): t\ge t_{\mathrm{ret}}(\mathbf{r})\big\}
\end{equation}
expresses finite propagation speed \(c_0=1/\sqrt{\mu_0\varepsilon_0}\) on vacuum
regions \cite{stratton1941}. Figure~\ref{fig:ch2.causal-light-cone}
schematizes Eq.~\eqref{eq:ch2.causal-support-constraint}.
\begin{figure}[t]
\centering
\ChOneFigBlock{%
\begin{tikzpicture}[ch1fig, node distance=5mm and 7mm]
 \draw[ch1flow] (0,0) -- (2.2,2.2);
 \draw[ch1flow] (0,0) -- (2.2,-2.2);
 \node[ch1oncanvas, right=2mm] at (2.2,2.2) {$t=t_{\mathrm{ret}}$};
 \node[ch1box] at (0,0) {Source event};
\end{tikzpicture}%
}
\caption{Causal support: field at \((\mathbf{r},t)\) depends on sources inside the
backward light cone; retarded integration uses \(t'\le t\) only.}
\label{fig:ch2.causal-light-cone}
\end{figure}
Energy balance in time domain imports the Poynting and storage laws of
Section~\ref{sec:ch01.1.4}. With volume energy density
\begin{equation}
\label{eq:ch2.time-energy-density}
u(\mathbf{r},t)
=
\frac{\varepsilon_0}{2}\|\E(\mathbf{r},t)\|^2
+\frac{\mu_0}{2}\|\H(\mathbf{r},t)\|^2,
\end{equation}
integration over a control volume \(V\subset\Omega\) yields
\begin{equation}
\label{eq:ch2.time-poynting-balance}
\frac{d}{dt}\int_V u\,dV
+\int_{\partial V}\mathbf{S}\cdot\hat{\mathbf{n}}\,dS
=
-\int_V \E\cdot\mathbf{j}_{\mathrm{imp}}\,dV,
\end{equation}
with \(\mathbf{S}=\E\times\H\) \cite{stratton1941}. Equation
\eqref{eq:ch2.time-poynting-balance} is cited; it supplies the physical
reading of Eq.~\eqref{eq:ch2.time-evolution-forced}: impressed work on the right
drives stored and exported energy on the left. Reactive exchange with boundaries
appears in the surface term and motivates the regime split of
Subsection~\ref{sec:ch224} \cite{harrington2001}. Harmonic operators follow from the substitution
\(\E(\mathbf{r},t)=\Re\{\underline{\E}(\mathbf{r})\eim\}\) with
\(\underline{\H}\) analogous, yielding
\begin{equation}
\label{eq:ch2.time-to-frequency-map}
\partial_t \longrightarrow -\jj\omega,
\qquad
\mathcal{A}\longrightarrow \mathcal{M}_\omega,
\qquad
\mathcal{S}_t\big|_{\omega}\longrightarrow \mathcal{S}_\omega,
\end{equation}
on admissible data of Eq.~\eqref{eq:ch2.admissibility-embedding}
\cite{stratton1941}. Equation~\eqref{eq:ch2.time-to-frequency-map} is symbolic;
Subsection~\ref{sec:ch223} states equivalence conditions under which frequency-domain
solvers implement causal time-domain maps. Exterior domains require outgoing
boundary conditions in time domain compatible with
Eq.~\eqref{eq:ch1.radiation-trace-condition}; those completions are sharpened in
Section~\ref{sec:ch243} \cite{stratton1941}.
Time-domain Maxwell-closed couples Eq.~\eqref{eq:ch1.admissible-space} at each time
slice to Eq.~\eqref{eq:ch2.charge-continuity-time}:
\begin{equation}
\label{eq:ch2.time-charge-continuous-coupling}
\mathbf{U}(t)\in\mathcal{F}_\Omega\ \text{a.e.\ }t
\ \wedge\
(\mathbf{j}_{\mathrm{imp}},\rho)\ \text{satisfy charge continuity}
\ \Longrightarrow\
\mathbf{U}\in\mathcal{D}(\mathcal{S}_{t}).
\end{equation}
Finite-difference time-domain schemes discretize
Eq.~\eqref{eq:ch2.time-maxwell-evolution} on staggered grids; stability requires
discrete analogues of Eq.~\eqref{eq:ch2.causal-support-constraint} and consistent
boundary enforcement of \(\mathcal{B}_\Gamma\) \cite{harrington2001}. Subsection
~\ref{sec:ch221} states the continuous targets those schemes approximate.
Subsection~\ref{sec:ch221} has posed the time-domain Maxwell system, the evolution
semigroup \(\mathcal{T}(t)\), and the retarded convolution form that enforces causal
support. Subsection~\ref{sec:ch222} carries the same physics into fixed-frequency
operators and Helmholtz symbols on the spaces of Section~\ref{sec:ch21}
\cite{stratton1941}.
